\begin{showBackmatter}%
%%%%%%%%%%%%%%%%%%%%%%%%%%%%%%%%%%%%%%%%
%%%%             Anhang             %%%%
%%%%%%%%%%%%%%%%%%%%%%%%%%%%%%%%%%%%%%%%
%\backmatter%
%
%% Falls man an dieser Stelle einen 
%% Seitenumbruch im Inhaltsverzeichnis
%% einfügen möchte:
%\addtocontents{toc}{\protect\newpage} 
%
%
%%%%%%%%%%%%%%%%%%%%%%%%%%%%%%%%%%%%%%%%
%%        Eigene Publikationen        %%
%%%%%%%%%%%%%%%%%%%%%%%%%%%%%%%%%%%%%%%%
\begin{showOwnPublications}%
%
%% Hier eigene Publikationen auflisten (aus der bib-Datei, die in Der Datei
%% .preambel/AllePfade.tex in \bibpathOwnPublications definiert wird).
\begin{refsegment}
%
%% Beim Erwähnen der Einträge in einer einleitenden Notiz:
%% keine Backreferenzen in der Liste erzeugen:
%\backtrackerfalse
%
%% Aufruf der Quellen ohne Verweis im Text geschieht mit \nocite{...}:
%%
%%  -  entweder nur die benötigten Quellen einzeln mit \nocite{X,Y,Z} aufrufen, z.B.
%\nocite{GrinbergOhrBeyerer-ITSC2009,TeutschGrinberg-RobustDetectionInWAMI-2016}
%%
%%  -  oder alle Quellen auf einmal mit \nocite{*} aufrufen und (sofern keine
%%     eigenständige, sondern eine gemeinsame Literaturdatenbank verwendet wird)
%%     ein Keyword (z.B. "ownpubl") als Filter verwenden
\nocite{*}
%
%% Es lässt sich bei Bedarf eine besondere Notiz hinzufügen:
%% Aktivierung erfolgt als Option von \printbibliography weiter unten
\defbibnote{OwnPublPrenote}{%
Dieser Abschnitt enthält ein vollständiges Verzeichnis der eigenen Veröffentlichungen.
%welche während der Promotionszeit entstanden sind.
Die Publikationen
%% Bei Referenzierung von Publikationen in diesem Abschnitt bitte 
%% statt dem üblichen \cite{...}-Makro das Makro
%% \citefield{...}[labelnumberwidth]{labelnumber} mit den beiden 
%% Optionen "labelnumberwidth" und "labelnumber" verwenden, damit
%% die Nummern und nicht die alphanumerischen Bezeichnungen zitiert werden:
\citePrefixed{GrinbergOhrBeyerer-ITSC2009,TeutschGrinberg-RobustDetectionInWAMI-2016}
adressieren das Thema XXX, während die Publikationen YYY sich mit dem Thema ZZZ beschäftigen. 
Die letzteren können nur im weiteren Sinne als relevant für die vorliegende Arbeit angesehen werden.
}
%
\appto{\bibsetup}{\RaggedRight} %Flattersatz, Silbentrennung zulassen
%\appto{\bibsetup}{\raggedright} %Flattersatz, keine Silbentrennung zulassen
\begin{refcontext}[sorting=ydnt, labelprefix=P]
\printbibliography[env=bibliographyNUM,%
				   segment=\therefsegment,
				   keyword=ownpubl,%
				   heading=bibintoc,%
				   %heading=subbibliography,%
				   %nottype=patent,%
				   prenote={OwnPublPrenote}, % Falls besondere Notiz (s.o.) erwünscht
				   resetnumbers=false,
				   title={\TransOwnPublications}]
\end{refcontext}
\end{refsegment}
\end{showOwnPublications}%
%
%
%%%%%%%%%%%%%%%%%%%%%%%%%%%%%%%%%%%%%%%%
%%           Eigene Patente           %%
%%%%%%%%%%%%%%%%%%%%%%%%%%%%%%%%%%%%%%%%
\begin{showOwnPatents}%
%
%% Hier eigene Patente auflisten (aus der bib-Datei, die in Der Datei
%% .preambel/AllePfade.tex in \bibpathOwnPatents definiert wird).
\begin{refsegment}
%
%% Beim Erwähnen der Einträge in einer einleitenden Notiz:
%% keine Backreferenzen in der Liste erzeugen:
%\backtrackerfalse
%
%% Aufruf der Quellen ohne Verweis im Text geschieht mit \nocite{...}:
%%
%%  -  entweder nur die benötigten Quellen einzeln mit \nocite{X,Y,Z} aufrufen, z.B.
%\nocite{PatentValeoUltraschall2016DE,PatentValeoUltraschall2016WO}
%%
%%  -  oder alle Quellen auf einmal mit \nocite{*} aufrufen und (sofern keine
%%     eigenständige, sondern eine gemeinsame Literaturdatenbank verwendet wird)
%%     ein Keyword (z.B. "ownpatent") als Filter verwenden
\nocite{*}
%
%% Es lässt sich bei Bedarf eine besondere Notiz hinzufügen:
%% Aktivierung erfolgt als Option von \printbibliography weiter unten
\defbibnote{OwnPatentPrenote}{%
Dieser Abschnitt enthält Patente, welche während meiner Promotionszeit entstanden sind.
Die Patente
%% Präfix und Nummer gemeinsam ausgeben, z.B. [Pat1]:
\citePrefixed{PatentValeoUltraschall2016DE,PatentValeoUltraschall2016WO}
entstanden in Zusammenarbeit mit...
}
%
\appto{\bibsetup}{\RaggedRight} %Flattersatz, Silbentrennung zulassen
%\appto{\bibsetup}{\raggedright} %Flattersatz, keine Silbentrennung zulassen
\begin{refcontext}[sorting=ydnt, labelprefix=Pat]
\printbibliography[env=bibliographyNUM,%
									 segment=\therefsegment,
									 keyword=ownpatent,%
									 %heading=subbibliography,%
									 heading=bibintoc,%
									 type=patent,%
									 prenote=OwnPatentPrenote, % Falls besondere Notiz (s.o.) erwünscht
									 resetnumbers=true,
									 title={\TransOwnPatents}]
\end{refcontext}
\end{refsegment}
\end{showOwnPatents}%
%
%
%%%%%%%%%%%%%%%%%%%%%%%%%%%%%%%%%%%%%%%%
%%   Betreute studentische Arbeiten   %%
%%%%%%%%%%%%%%%%%%%%%%%%%%%%%%%%%%%%%%%%
\begin{showSupervisedTheses}%
%%
\begin{refsegment}
%
%% Aufruf der Quellen ohne Verweis im Text:
%
%% entweder nur die benötigten Quellen einzeln mit \nocite{X,Y,Z} aufrufen
%\nocite{%
%Mustermann-2016,
%Mustermann-2017}
%
%% oder alle Quellen auf einmal mit \nocite{*} aufrufen
%% und ein Keyword (z.B. "supervised") als Filter verwenden
\nocite{*}
%%
\defbibnote{SupervThesisPrenote}{%
%% Bei Referenzierung von Publikationen in diesem Abschnitt bitte 
%% statt dem üblichen \cite{...}-Makro das Makro
%% \citefield{...}[labelnumberwidth]{labelnumber} verwenden, damit
%% die Nummern und nicht die alphanumerischen Bezeichnungen zitiert werden.
Dieser Abschnitt enthält die während der Promotionszeit betreuten studentischen Arbeiten.
}
%%
\appto{\bibsetup}{\RaggedRight} %Flattersatz, Silbentrennung zulassen
%\appto{\bibsetup}{\raggedright} %Flattersatz, keine Silbentrennung zulassen
\begin{refcontext}[sorting=ydnt, labelprefix=T]
\printbibliography[env=thesislist,%bibliographyNUM,%
									 segment=\therefsegment,
									 keyword=supervised,%
									 %heading=subbibliography,%
									 %type=thesis,%
									 heading=bibintoc,%
									 prenote=SupervThesisPrenote, % Falls besondere Notiz (s.o.) erwünscht
									 resetnumbers=true,
									 title={\TransSupervisedTheses}]
\end{refcontext}
\end{refsegment}
\end{showSupervisedTheses}%
%
%
%%%%%%%%%%%%%%%%%%%%%%%%%%%%%%%%%%%%%%%%
%%%%          Bibliographie         %%%%
%%% (Achtung! Bibliographie darf im  %%%
%%%  deutschen nicht in den Anhang!) %%%
%%%%%%%%%%%%%%%%%%%%%%%%%%%%%%%%%%%%%%%%
%\newrefcontext[sorting=none]
\begin{showBibliography}%
\begin{refsegment}
%\backtrackertrue
%% mit BibLaTeX:
\appto{\bibsetup}{\RaggedRight} %Flattersatz, Silbentrennung zulassen
%\appto{\bibsetup}{\raggedright} %Flattersatz, keine Silbentrennung zulassen
\begin{refcontext}[sorting=none]
\printbibliography[env=bibliographyNUM,
					heading=bibintoc,
					category=cited,
					notkeyword=supervised,
					notkeyword=ownpubl,
					notkeyword=ownpatent,
					resetnumbers=true
					]
\end{refcontext}
\end{refsegment}
\end{showBibliography}%
%
%
%%%%%%%%%%%%%%%%%%%%%%%%%%%%%%%%%%%%%%%%
% Abbildungs- und Tabellenverzeichnis  %
%%%%%%%%%%%%%%%%%%%%%%%%%%%%%%%%%%%%%%%%
\begin{showListOfFigures}
\listoffigures
\cleardoublepage
\end{showListOfFigures}%
%
\begin{showListOfTables}
\listoftables
\cleardoublepage
\end{showListOfTables}%
%
\begin{showListOfListings}
\lstlistoflistings
\cleardoublepage
\end{showListOfListings}%
%
%
%%%%%%%%%%%%%%%%%%%%%%%%%%%%%%%%%%%%%%%%
%%%%  Alle Glossare etc. ausgeben:  %%%%
%%%%%%%%%%%%%%%%%%%%%%%%%%%%%%%%%%%%%%%%
%\begin{showGlossaries}%
%\IfDefined{printglossaries}{%
%	\printglossaries[style=mylongglossstyle]
%}
%\end{showGlossaries}%
%%%%%%%%%%%%%%%%%%%%%%%%%%%%%%%%%%%%%%%%
%
%
%
% Oder einzelne Glossartypen ausgeben:
%
%%%%%%%%%%%%%%%%%%%%%%%%%%%%%%%%%%%%%%%%
%%%%          Abkürzungen           %%%%
%%%%%%%%%%%%%%%%%%%%%%%%%%%%%%%%%%%%%%%%
\begin{showListOfAcronyms}%
\IfDefined{printglossary}{%
	\setlength{\LTleft}{0pt}% Einzug entfernen (Forderung des KSP-Verlages)
	%\setlength{\LTleft}{-5pt}% Bei Verwendung von style=long-booktabs
	%\printacronyms
	%% Falls das Abkürzungsverzeichnis eigene Breiten der Namen-Spalten haben soll:
	%\glssetwidest[0]{Haupteintrag}% <--- Hier längstes Wort unter den Glossar-Einträgen angeben
	%\glssetwidest[1]{Untereintraglänge}% <--- Hier längstes Wort unter den Glossar-Untereinträgen der 1 Ebene
	%\glssetwidest[2]{Untereintraglänge}% <--- Hier längstes Wort unter den Glossar-Untereinträgen der 2 Ebene
	\printglossary[type=\acronymtype,style=alttree]
}%
\cleardoublepage
\end{showListOfAcronyms}%
%
%%%%%%%%%%%%%%%%%%%%%%%%%%%%%%%%%%%%%%%%
%%%%      Begriffsverzeichnis       %%%%
%%%%%%%%%%%%%%%%%%%%%%%%%%%%%%%%%%%%%%%%
\begin{showNomenclature}%
\IfDefined{printnomenclature}{%
	\setlength{\LTleft}{0pt}% Einzug entfernen (Forderung des KSP-Verlages)
	%\setlength{\LTleft}{-5pt}% Bei Verwendung von style=long-booktabs
	\printnomenclature
}%
\cleardoublepage
\end{showNomenclature}%
%
%%%%%%%%%%%%%%%%%%%%%%%%%%%%%%%%%%%%%%%%
%%%%            Glossar             %%%%
%%%%%%%%%%%%%%%%%%%%%%%%%%%%%%%%%%%%%%%%
\begin{showGlossary}%
\IfDefined{printglossary}{%
	\setlength{\LTleft}{0pt}% Einzug entfernen (Forderung des KSP-Verlages)
	%\setlength{\LTleft}{-5pt}% Bei Verwendung von style=long-booktabs
	%\printglossary
	%\printglossary[style=super]
	%\printglossary[style=longheader]
	%\printglossary[style=long]
	%\printglossary[style=mylongglossstyle]
	%% Falls der Glossar eigene Breiten der Namen-Spalten haben soll:
	%\glssetwidest[0]{Haupteintrag}% <--- Hier längstes Wort unter den Glossar-Einträgen angeben
	%\glssetwidest[1]{Untereintraglänge}% <--- Hier längstes Wort unter den Glossar-Untereinträgen der 1 Ebene
	%\glssetwidest[2]{Untereintraglänge}% <--- Hier längstes Wort unter den Glossar-Untereinträgen der 2 Ebene
	\printglossary[style=alttree]
}%
\cleardoublepage
\end{showGlossary}%
%
%
%
%
%
%%%%%%%%%%%%%%%%%%%%%%%%%%%%%%%%%%%%%%%%
%%%%             Anhang             %%%%
%%%%%%%%%%%%%%%%%%%%%%%%%%%%%%%%%%%%%%%%
\begin{showAppendix}
\appendix
% Platziere einen Verweis auf 'Anhang'
% ins Inhaltsverzeichnis
\phantomsection
\addcontentsline{toc}{part}{\AppendixName}
%
%%%%%%%%%%%%%%%%%%%%%%%%%%%%%%%%%%%%%
\chapter[Herleitungen]{Herleitungen}%
\label{chap:AppendixDerivations}
%%%%%%%%%%%%%%%%%%%%%%%%%%%%%%%%%%%%%
%
Hier kommen die Herleitungen rein.
%
% Reset the PDF bookmark level
\bookmarksetup{startatroot}
\cleardoublepage
\end{showAppendix}%
%
%
%
%%%%%%%%%%%%%%%%%%%%%%%%%%%%%%%%%%%%%%%%
%%%%     Stichwortverzeichnis       %%%%
%%%%%%%%%%%%%%%%%%%%%%%%%%%%%%%%%%%%%%%%
\begin{showIndex}%
% Platziere einen Verweis auf das Stichwortverzeichnis
% ins Inhaltsverzeichnis
\phantomsection
\IfDefined{printindex}{%
	\addcontentsline{toc}{chapter}{\indexname}%
	\setlength{\LTleft}{-5pt}% Einzug entfernen (Forderung des KSP-Verlages)
	%\renewcommand{\indexname}{\TransIndexName}%
	\printindex
	\cleardoublepage
}%
\end{showIndex}%
%
%
%
%%%%%%%%%%%%%%%%%%%%%%%%%%%%%%%%%%%%%%%%
%%%%         Liste der TODOs        %%%%
%%%%%%%%%%%%%%%%%%%%%%%%%%%%%%%%%%%%%%%%
\begin{showTODOs}
\phantomsection
% Zum Inhaltsverzeichnis hinzufügen
\todototoc
\listoftodos%
\hypertarget{TODO-List}{}%
\end{showTODOs}%
%
\end{showBackmatter}%